% !Mode:: "TeX:UTF-8"

\chapter{结\quad 论}

\section{全文总结}

本文构建了一个面向大语言模型空间推理行为的数据挖掘框架。项目通过多层级 Prompt、无关噪声扰动和可解释行为特征，能够从准确率、逻辑一致性、噪声敏感性和失效模式等角度分析模型能力边界。当前版本已经形成完整的代码流程和报告结构，后续可进一步补充真实模型多轮采样结果，并对关键发现进行统计检验。

\section{未来展望}

后续工作可以从三方面展开：第一，扩大模型和重复采样数量，提高结论稳定性；第二，引入更严格的几何基准或人工复核机制，提升参考标签可信度；第三，将该方法扩展到电梯通行、门洞转向、房间路径规划等更多空间任务，以验证模型空间建模能力的泛化边界。
